\chapter{Conclusiones del avance}
La ampliación del corpus permite formular un problema más preciso. La literatura reunida ya contiene alternativas que combinan reducción, clasificación y explicabilidad, además de trabajos que miden recursos del modelo. El seminario debe reconocer esos avances y concentrarse en una pregunta comprobable: cuánto cuesta explicar y qué beneficio adicional ofrece integrar selección de atributos, destilación y asignación adaptativa de cómputo.

La arquitectura multinivel se mantiene, pero su evaluación adquiere un papel central. TreeSHAP completo y el estudiante entrenado directamente serán referencias necesarias para atribuir los efectos observados. La confianza del modelo se tratará como una señal que necesita calibración y auditoría. La reducción del número de variables mostradas no se confundirá con ahorro de cálculo, y las explicaciones del estudiante no se presentarán como explicaciones del profesor.

El aporte de esta etapa es documental y metodológico. Se ha definido qué se comparará, cómo se separarán los costos y qué controles permitirán evaluar rendimiento y calidad técnica. Quedan pendientes la depuración bibliográfica de los antecedentes antiguos, la auditoría de datos y la implementación experimental. Por tanto, todavía no corresponde afirmar que el sistema reduce latencia, conserva exactitud o mejora decisiones humanas.

La siguiente etapa deberá convertir este diseño en evidencia reproducible. Si el enfoque propuesto aporta una ventaja, será posible describir sus condiciones y límites. Si una solución más sencilla ofrece un equilibrio equivalente o mejor, la investigación también habrá producido una respuesta útil para diseñar clasificación de malware explicable con un costo justificado.
